\documentclass[12pt]{scrartcl}
\usepackage[sexy]{evan}
\newcommand{\R}{\mathbb{R}}
\newcommand{\N}{\mathbb{N}}
\newcommand{\Q}{\mathbb{Q}}
\newcommand{\Z}{\mathbb{Z}}
\usepackage{relsize}
\begin{document}
\title{Syllabus}
\author{Jiawen Huang}
\pagenumbering{gobble} % Remove page numbers
\vspace{0.1cm}
\begin{center}
    \LARGE \textbf{Syllabus}
\end{center}

\section{Day 1}
\begin{enumerate}
    \item Naive Set Theory: disjoint union
    \item Equivalence relations
    \item Canonical decomposition of a map
    \item Definition of category
    \item Isomorphisms
    \item Initial and terminal objects
    \item A commutative diagram
    \item No HW. 
\end{enumerate}
\section{Day 2}
\begin{enumerate}
    \item What is a group?
    \item Examples of groups 
    \item No HW. 
\end{enumerate}

\section{Day 3}
\begin{enumerate}
    \item Order of an element
    \item Examples of order of elements
    \item Property: if $g^n=e$, then order of $g$ divides $n$.
    \item symmetric groups (automorphism of set of 1 to n, n!, $S_3$ )
    \item dihedral groups $D_{2n}$ (group element is operation)
    \item cyclic groups ($\Z/n\Z$, who is the generator, order n)
    \item units group $(\Z/n\Z)^\times$, is there a generator? that is, is it cyclic? (this is very hard)
    \item 
    \item HW: reading chapter 2 sections 1 and 2. Do 1.3, 1.4, 1.5, 1.8, 1.10, 1.11. Do 2.2, 2.5, 2.6, 2.7, 2.11, 2.12, 2.13, 2.15, 2.17
\end{enumerate}

\section{Day 4}
\begin{enumerate}
    \item Group homomorphisms. $\phi(a*b) = \phi(a)\times \phi(b)$. Show $\phi(e)=e'$, $\phi(a^{-1})=\phi(a)^{-1}$ 
    \item Example: $e$ as initial and final. $\ZZ$ to $\ZZ/n\ZZ$, $\ZZ$ to $G$ by $g^n$. $(\RR,+)$ to $(\RR^\times, \times)$, $S_3$ to $D_6$, $S_3$ to $C_2$. $\det(M)$ as homomorphism from $GL_n(\R)$ to $\R^\times$. $\CC^\times$ to $\R^\times$ by $z\mapsto |z|$. 
    \item Category of groups. 
    \item Category of abelian groups. 
    \item Initial and final objects in groups: trivial group
    \item Product groups: product is a universal object in some sense. What is product in sets? What is product in general. Product and coproduct in groups and abelian groups. 
    \item Some homomorphisms from and to product groups.
    \item Group homomorphisms and order of elements. 
    \item Isomorphisms in general
    \item groups: bijective homomorphisms (nontrivial, deserve a proof): $S_3\cong D_6$, $\Z/6\Z \not\cong D_6$, $(\Z/7\Z)^\times \cong \Z/6\Z$, $C_2\times C_3 \cong C_6$. 
    \item homeomorphisms of topological spaces is not bijective morphism. 
    \item When not isomorphic? Key property, isomorphism preserves order of an element. Look at order of elements.
    \item homomorphisms of abelian groups form a group, in fact, a ring. 
    \item automorphism of groups form a group.
    \item HW: reading chapter 2,  section 3, 4. Do 3.2, 3.4. Do 4.1, 4.3, 4.7, 4.8, 4.9, 4.10, 4.13, 4.14, 4.15, 4.17. 
\end{enumerate}

\section{Day 5}
\begin{enumerate}
    \item Universal properties in category theory: Given a catgeory, a natural question is whether there are initial and final (not always exists). We have seen that initial and final objects are unique up to isomorphisms which is very nice. Inital and final objects are "extreme" objects in the category without posing any extra conditions. We can also ask for "extreme" objects with some extra conditions, and this is what universal properties are about: the way to do this is to construct a new category, very much similar to the old one but posing the condition we need. And we find the initial object in this new category. For example, we want to find an "extreme" object $P$ with morphisms to $A$ and $B$ ($A$ and $B$ are fixed) We can ask for the "extreme" object with some extra conditions, for example, we want the morphism from $P$ to $A$ and $B$ to be homomorphisms. This is what we call product of $A$ and $B$.  
    \item Products in groups, sets, abelian groups and topological spaces.
    \item Coproducts in sets, abelian groups and topological spaces. 
    \item HW: reading (chapter 1, section 5). Do all exericses form ch1 section 5. 
\end{enumerate}
\section{Day 6}

\begin{enumerate}
    \item Products in topological spaces.
    \item Coproducts in sets, abelian groups and topological spaces. 
    \item Free groups and free abelian groups. 
    \item Universal property of free groups, free groups is universal so unique up to isomorphisms.
    \item Construction of free groups (maybe skip)
    \item Construction of free abelain groups. 
    \item Chapter 2 section 5 Do 5.1, 5.2, 5.3, 5.4, 5.5, 5.9, 5.10 in chapter 2 section 5. Exercise: find the free objects in the category of topological spaces. 
\end{enumerate}
\section{Day 7}
\begin{enumerate}
    \item Definition of subgroups in terms of category theory: a subgroup of $G$ is a group $H$ with an injective homomorphism to $G$.
    \item Standard definition of subgroup: a subset of $G$ closed under group operation and inverse.
    \item $H$ is a subgroup of $G$ if and only if for any $a,b\in H$, $ab^{-1}\in H$. The most useful criterion for checking whether a subset is a subgroup.
    \item $\phi\inv(H)$ is a subgroup of $G$ if $H$ is a subgroup of $G'$ given a homomorphism $\phi:G\to G'$.
    \item Kernel and image of a homomorphism are subgroups.
    \item Example: subgroup of $S_3$,  $\ZZ/n\ZZ$, $\CC^\times$, $C^0([0,1])$.
    \item Subgroup generated by a subset. Finitely generated groups: there exists $F(\{a_1, \ldots, a_n\})\surjto H$.
    \item Subgroup of $\ZZ$ is nice. 
    \item Injective is just kernel is trivial.
    \item HW: read section 6. Do 6.4, 6.5, 6.6, 6.7, 6.8, 6.9, 6.12, 6.14, 6.15, 6.16. 
\end{enumerate}
\end{document}